\chapter{Control Térmico}

\section{Mecanismos de Control Térmico LEO}
El control térmico del nanosatélite INTISAT es transversal; protege al hardware de sobrevivir a los extremos ciclos orbitales intermedios entre iluminación solar directa severa ($+100\ ^\circ\text{C}$) y profundos eclipses de umbra por la Tierra ($-40\ ^\circ\text{C}$). Se divide en dos filosofías conectadas al hardware de la nave:

\subsection{Control Pasivo (Recubrimientos Ópticos y Masa)}
Delegado inherentemente a la construcción del \textit{Chasis (Estructura)}:
\begin{itemize}
    \item \textbf{Radiadores:} zonas de aluminio no anodizado (superficies desnudas o pulidas que se dejan al final del ensamblaje geométrico) por donde se disipa el flujo de calor sobrante de la batería hacia el vacío.
    \item \textbf{Aislamiento y Recubrimientos Ópticos:} Empleo de cintas Kapton o aluminizadas para blindar internamente las caras del PCB sensibles al puente frío exterior. 
\end{itemize}

\subsection{Control Activo (Integración EPS y OBC)}
Para proteger permanentemente las celdas de las baterías contra las bajas temperaturas en período de eclipse (que causarían congelamiento del electrolito y pérdida de capacidad):
\begin{itemize}
    \item \textbf{Calentadores Flexibles (Heaters):} Parches resistivos adheridos directamente al *Battery Module*. Su activación la orquesta el sistema de control EPS mediante la realimentación indirecta del BMS local y los 6 sensores RTD de temperatura de la arquitectura central del OBC.
\end{itemize}

\section{Campaña de Pruebas Térmicas (AIT)}
El comportamiento se simula en programas CAD pero se rectifica únicamente de manera empírica.

\subsection{Ciclos Térmicos (Thermal Cycling)}
Se somete al satélite acoplado a pruebas dinámicas en cámara, ciclando la temperatura (por ejemplo, entre $-15\ ^\circ\text{C}$ y $+60\ ^\circ\text{C}$) reiteradas veces (ej. 10 ciclos). El OBC debe continuar transmitiendo a través de la interfaz umbilical para detectar termofatiga en soldaduras perimetrales críticas.

\subsection{Prueba de Vacío Térmico (Bake-Out LEO)}
En sincronía con la agencia inyectora, el hardware es horneado al alto vacío extremo durante semanas a $+60\ ^\circ\text{C}$ promoviendo la máxima desgasificación (\textit{outgassing}). Esta depuración es obligatoria para certificar que el satélite prevenga la contaminación de la óptica de los *Payloads microscópicos* e instrumentos vecinos en el cohete.
